\chapter{Implémentation logicielle}
\label{chap:implementation}

\section{Choix techniques}

\begin{table}[H]
\centering
\caption{Choix techniques et justifications}
\label{tab:choix-techniques}
\begin{tabular}{@{}llp{7cm}@{}}
\toprule
\textbf{Besoin} & \textbf{Retenu} & \textbf{Justification} \\
\midrule
Langage & Python 3.11 & écosystème scientifique, lisibilité, compétence disponible \\
Modélisation & PuLP & syntaxe proche de la notation mathématique \\
Solveur & CBC & libre, sans licence : condition de pérennité chez OCP \\
Tests & pytest & standard, paramétrage et fixtures simples \\
Restitution & openpyxl & format imposé par le service \\
\bottomrule
\end{tabular}
\end{table}

Le problème comptant \num{29} variables binaires et \num{138} contraintes, la puissance d'un
solveur commercial est inutile. PuLP permettrait au demeurant de basculer sur Gurobi ou
CPLEX en changeant une seule ligne, si la taille du problème venait à croître.

\section{Architecture}

\begin{figure}[H]
\centering
\resizebox{\textwidth}{!}{%
\begin{tikzpicture}[
  node distance=0.75cm,
  mod/.style={rectangle, draw=ocpbleu, thick, fill=white, rounded corners=2pt,
              text width=4.6cm, align=center, minimum height=0.85cm, font=\small},
  ann/.style={font=\scriptsize\itshape, text=grisfonce, align=left, text width=5.4cm},
  f/.style={-{Stealth[length=2.5mm]}, thick, draw=ocpbleu},
]
\node[mod] (sc) {\texttt{scenario.py}};
\node[ann, right=0.5cm of sc] (asc) {validation stricte : un scénario incohérent échoue tout de suite};
\node[mod, below=of sc] (pre) {\texttt{preprocessing.py}};
\node[ann, right=0.5cm of pre] (apre) {prétraitement déterministe : $\Pi_m$, $G_m$, $R_{k,a}$};
\node[mod, below=of pre] (mod) {\texttt{model/}};
\node[ann, right=0.5cm of mod] (amod) {MILP : 170 variables, 138 contraintes};
\node[mod, below=of mod] (sol) {\texttt{solver.py}};
\node[ann, right=0.5cm of sol] (asol) {trois passes lexicographiques (CBC)};
\node[mod, below=of sol] (res) {\texttt{results.py}};
\node[ann, right=0.5cm of res] (ares) {vue métier de la solution};
\node[mod, below=1.3cm of res] (val) {\texttt{validation.py}};
\node[mod, left=0.7cm of val] (ana) {\texttt{analysis.py}};
\node[mod, right=0.7cm of val] (rep) {\texttt{reporting/}};

\draw[f] (sc) -- (pre);
\draw[f] (pre) -- (mod);
\draw[f] (mod) -- (sol);
\draw[f] (sol) -- (res);
\draw[f] (res) -- (val);
\draw[f] (res.south west) to[out=-120,in=90] (ana.north);
\draw[f] (res.south east) to[out=-60,in=90] (rep.north);
\end{tikzpicture}%
}
\caption{Chaîne de traitement du logiciel}
\label{fig:architecture}
\end{figure}

Le paquet compte treize modules répartis comme suit : \texttt{constants}, \texttt{units},
\texttt{profiles}, \texttt{scenario} et \texttt{preprocessing} pour le socle ; le
sous-paquet \texttt{model/} (variables, contraintes, objectif) ; \texttt{solver},
\texttt{results}, \texttt{validation}, \texttt{analysis} et \texttt{cli} ; le sous-paquet
\texttt{reporting/} pour les sorties Excel et JSON.

\section{Trois décisions d'architecture}

\subsection{Le prétraitement est séparé de l'optimisation}

Une grande partie du système n'est pas optimisable. Ces grandeurs sont calculées une fois
pour toutes \emph{avant} de construire le modèle.

\begin{encadre}[title={Pourquoi c'est important}]
En faire des variables de décision donnerait un modèle plus gros, plus lent, et surtout
\textbf{faux} : il produirait des plans que l'usine ne peut pas exécuter. C'est l'erreur la
plus facile à commettre sur ce sujet.
\end{encadre}

\subsection{Une fonction par famille de contraintes, nommée d'après sa référence}

Les fonctions s'appellent \texttt{c07\_bilan\_acide\_29\_std},
\texttt{c11\_capacite\_decanteurs}, et ainsi de suite. La correspondance entre le
chapitre~\ref{chap:modelisation} et le code se vérifie ligne à ligne : une contrainte
modifiée dans le modèle se retrouve immédiatement dans le code, et réciproquement.

\subsection{Le validateur ne relit jamais le modèle}

Le module \texttt{validation.py} recalcule les bilans \textbf{à partir des équations
physiques}, et non des contraintes posées dans \texttt{model/}.

\begin{encadre}[title={Le principe de double vérification}]
Si la formulation contient une erreur, le solveur renverra une solution parfaitement
« optimale » --- pour un mauvais problème. Rejouer les contraintes du modèle ne détecterait
rien. Seul un contrôle écrit séparément peut mettre l'erreur en évidence.
\end{encadre}

\begin{resultat}[title={Ce principe a payé dès le premier passage}]
Le validateur a immédiatement signalé un écart de \SI{25,7}{\tonne} sur le bilan d'IR11.
L'origine : dans l'extraction des résultats, les livraisons de CoC et de DEC\_CL vers un
\emph{même} consommateur s'écrasaient mutuellement au lieu de se cumuler --- une
compréhension de dictionnaire écrivant deux fois la même clé.

Le modèle, lui, était correct. C'est la couche de restitution qui était fautive. Un contrôle
qui aurait relu les contraintes du modèle n'aurait rien vu.
\end{resultat}

\section{Points délicats de l'implémentation}

\subsection{Le piège de portée en Python}

Dans une fonction imbriquée qui ajoute des contraintes, écrire \texttt{prob += ...}
rebinderait \texttt{prob} en variable \emph{locale} de la fonction interne, provoquant une
\texttt{UnboundLocalError}. D'où l'usage de \texttt{prob.addConstraint(...)}.

\subsection{Les tolérances de la résolution lexicographique}

Figer un niveau à l'exact (\texttt{f1 == f1*}) rend la passe suivante infaisable : le
solveur ne retrouve pas au bit près la valeur qu'il vient d'annoncer. On fige donc avec une
tolérance relative de \num{e-6}, négligeable devant des grandeurs en milliers de tonnes.

Ces contraintes de figeage sont \textbf{retirées en sortie}, y compris en cas d'erreur. Sans
ce nettoyage, la résolution laisserait le modèle dans un état différent de celui qu'elle a
reçu : une seconde résolution du même objet accumulerait les contraintes, et le décompte
publié dans ce rapport deviendrait faux. Le défaut a d'ailleurs été mis au jour par le
fichier de tests décrit en section~\ref{sec:tests-valeurs}.

\subsection{Un mode paramétrable plutôt qu'une bifurcation du code}

Le mode « cocristallisation décidable » (décision~D-11) aurait pu donner lieu à deux
variantes du modèle. On a préféré quatre fonctions d'aide --- \texttt{\_pi\_coc},
\texttt{\_coc\_produit}, \texttt{\_boue\_coc}, \texttt{\_coc\_total} --- qui renvoient soit
le paramètre, soit l'expression affine, selon que les binaires existent ou non :

\begin{lstlisting}[language=Python, caption={Bascule entre paramètre et variable}, label={lst:pi-coc}]
def _pi_coc(v, params, ligne54):
    if not v.z:                       # mode subi : c'est un parametre
        return params.production_54_coc[ligne54]
    return pulp.lpSum(                # mode decidable : expression affine
        params.production_echelon[e] * v.z[e]
        for e in params.echelons_coc_candidats.get(ligne54, ())
    )
\end{lstlisting}

Les fonctions de contraintes ignorent donc complètement le mode. C'est ce qui garantit
qu'aucune des deux variantes ne peut diverger de l'autre par inadvertance.

\subsection{Un écart de 0,06\,\% sur le coefficient de conversion}

Le dossier publie $\mathrm{vol}_{29}(h) = 174{,}11 \cdot h \cdot 0{,}33$, alors que la
physique déclarée donne $0{,}26 \times 1{,}270 = 0{,}3302$. Le code retient le coefficient
publié --- c'est la formule réellement utilisée sur le site --- et un test vérifie en
permanence qu'il reste cohérent avec la physique.

\section{Stratégie de test}

\num{162} tests s'exécutent en environ \SI{3}{\second}. Au-delà du décompte, quatre familles
méritent d'être distinguées.

\subsection{Les tests qui verrouillent une anomalie}

Présentés au chapitre~\ref{chap:analyse-specs}. Un constat documenté mais non testé finit
par être oublié.

\subsection{Les tests de contre-épreuve}

Ils démontrent qu'une contrainte \emph{sert à quelque chose} :

\begin{lstlisting}[language=Python, caption={Contre-épreuve de la contrainte de qualité d'IR11}, label={lst:contre-epreuve}]
def test_sans_contrainte_de_qualite_aucun_dec_cl(scenario, profils):
    """Avec alpha = 0, l'optimiseur cesse completement de produire du DEC_CL."""
    sans_alpha = dataclasses.replace(scenario, alpha_dec_cl=0.0)
    params = calculer_parametres(sans_alpha, profils)
    modele = construire_modele(sans_alpha, params)
    resultat = extraire(modele, resoudre(modele))

    dec_cl = sum(r.dec_cl_produit for r in resultat.lignes_54.values())
    assert dec_cl == pytest.approx(0.0, abs=1e-4)
\end{lstlisting}

Ce test \emph{prouve} la nécessité de la contrainte \Ctr{15} : sans elle, la règle métier
serait effectivement violée. L'argument théorique du
chapitre~\ref{chap:modelisation} est ainsi confirmé expérimentalement.

\subsection{Les tests qui corrompent la solution}

Un validateur qui répondrait toujours « conforme » passerait tous les tests positifs sans
rien garantir. On lui soumet donc des solutions délibérément fausses --- treize corruptions
différentes : bilan faussé, demande non servie, capacité dépassée, niveau de décadmiation
invalide, stock négatif, transfert sur une liaison inexistante --- et l'on vérifie qu'il les
rejette toutes.

\subsection{Les tests qui verrouillent les valeurs publiées}
\label{sec:tests-valeurs}

Ce projet cite des valeurs numériques dans une dizaine de documents et dans le présent
rapport. Une modification du modèle rendrait silencieusement fausse une partie de cette
documentation --- et une documentation fausse est pire qu'une documentation absente, car on
lui fait confiance.

Un fichier de tests dédié verrouille donc chaque chiffre publié : taille du modèle, stocks
initiaux, bornes de cuve, production d'acide 54, demande consolidée, valeurs des trois
niveaux, plan optimal, violations, tableaux de sensibilité. Chaque assertion porte en
commentaire l'endroit où la valeur est publiée.

\begin{encadre}
Ce filet a immédiatement servi : il a révélé que la résolution laissait dans le modèle ses
deux contraintes de figeage lexicographique, ce qui portait le décompte publié de 138 à 140
et rendait l'objet non réutilisable. Le correctif est décrit plus haut.
\end{encadre}

\section{Sorties produites}

Le programme produit une synthèse console, le classeur Excel à cinq feuilles imposé par le
service, et un export JSON. Ce dernier n'est pas un doublon : il permet de rejouer, comparer
et archiver les résultats sans passer par un tableur, ce qui est indispensable aux analyses
de sensibilité et aux tests de non-régression.

\begin{lstlisting}[language=bash, caption={Utilisation}, label={lst:usage}]
pip install -e ".[dev]"

python -m ocp_optim                        # resout et affiche la synthese
python -m ocp_optim --sortie resultats/    # rapports Excel et JSON
python -m ocp_optim --sensibilite          # analyses de sensibilite

pytest                                     # 162 tests, environ 1,2 s
\end{lstlisting}

Le code de retour vaut \num{0} si la validation est conforme et \num{1} sinon : le programme
est donc utilisable dans une chaîne automatisée.
